\documentclass[10pt,letterpaper]{article}

\usepackage[utf8]{inputenc}
\usepackage[T1]{fontenc}
\usepackage[spanish,es-nodecimaldot]{babel}
\usepackage[margin=0.48in]{geometry}
\usepackage{xcolor}
\usepackage{enumitem}
\usepackage{titlesec}
\usepackage{hyperref}

\hypersetup{
    colorlinks=false,
    pdfborder={0 0 0}
}

\pagestyle{empty}
\setlength{\parindent}{0pt}
\setlength{\parskip}{1pt}

\definecolor{headerblue}{RGB}{33, 47, 61}
\definecolor{accent}{RGB}{41, 128, 185}
\definecolor{graytext}{RGB}{85, 85, 85}

\titleformat{\section}
  {\normalsize\bfseries\color{headerblue}}
  {}
  {0em}
  {}
  [\vspace{-2pt}\color{accent}\titlerule]
\titlespacing*{\section}{0pt}{4pt}{2pt}

\begin{document}

\begin{center}
    {\Large\bfseries\color{headerblue} GABRIEL QUEZADA AVILA}\\[2pt]
    {\normalsize\bfseries\color{accent} Ingeniero en Informática}\\[2pt]
    {\footnotesize\color{graytext}
    Ciudad de México \quad|\quad 55 8719-5907 \quad|\quad gabrielqva.10@gmail.com \quad|\quad github.com/quezarbelia
    }
\end{center}

\vspace{-6pt}

\section{Perfil Profesional}
Ingeniero en Informática enfocado en Análisis de Datos, Business Intelligence y Ciencia de Datos. Experiencia en el sector público gestionando flujos de datos operativos y fiscales, diseño de bases de datos relacionales, automatización de procesos mediante scripts y desarrollo de tableros ejecutivos. Orientado a transformar requerimientos funcionales en soluciones analíticas estructuradas, gobernanza del dato, comunicación técnica clara y adopción de tecnologías en la nube.

\section{Experiencia Profesional}

\textbf{Analista de Datos y Operaciones} \hfill Agosto 2025 -- Actualidad\\
\textit{CONTRIBUTEL -- Secretaría de Administración y Finanzas de la CDMX (SAFCDMX)}
\begin{itemize}[leftmargin=12pt, itemsep=0pt, topsep=1pt, parsep=0pt]
    \item Estructuración, limpieza y depuración de flujos de datos operativos y fiscales para el control de la gestión institucional.
    \item Automatización y mantenimiento de reportes periódicos mediante consultas complejas en SQL y macros con VBA en MS Access.
    \item Generación de informes analíticos y tableros con métricas de rendimiento orientados a la toma de decisiones directivas.
    \item Asesoría especializada y resolución de incidencias normativas a contribuyentes, traduciendo procesos técnicos a lenguaje claro.
\end{itemize}

\textbf{Consultor de Sistemas y Datos (Proyectos Independientes)} \hfill Enero 2025 -- Actualidad\\
\textit{Desarrollo Autónomo y Arquitectura de Información}
\begin{itemize}[leftmargin=12pt, itemsep=0pt, topsep=1pt, parsep=0pt]
    \item Modelado e implementación de bases de datos relacionales (PostgreSQL, MySQL) con consultas optimizadas y triggers.
    \item Desarrollo de scripts en Python para la extracción, transformación y carga (ETL) de fuentes de datos heterogéneas.
    \item Diseño de flujos funcionales y prototipos interactivos en Figma integrados con bases de datos y servicios backend.
\end{itemize}

\section{Educación}

\textbf{Ingeniería en Informática} (Egresado) \hfill UPIICSA -- IPN\\
UPIICSA -- Instituto Politécnico Nacional (IPN) \hfill Ciudad de México

\textbf{Técnico en Informática} \hfill 2021 -- 2024\\
CECyT 5 ``Benito Juárez'' -- Instituto Politécnico Nacional (IPN) \hfill Ciudad de México

\section{Habilidades Técnicas}

\textbf{Análisis de Datos y BI:} Power BI (DAX, Power Query), Tableau Server / Desktop, Microsoft Fabric, Microsoft Excel Avanzado (Tablas dinámicas, modelado y Solver).\\
\textbf{Bases de Datos y DBA:} Oracle Database (administración básica/DBA, SQL, PL/SQL), Microsoft SQL Server, PostgreSQL, MySQL, SQLite, MS Access.\\
\textbf{Lenguajes de Programación:} Python, SQL, Java, C++, VBA.\\
\textbf{Cloud y Entornos (Nivel Práctico):} Google Cloud Platform (BigQuery), AWS (S3, RDS, EC2), Supabase, Linux (Ubuntu, Debian), Git, GitHub.\\
\textbf{Modelado y Diseño:} Diagramas Entidad-Relación (DER), arquitectura de datos, Figma.

\section{Certificaciones y Formación Continua}
\begin{itemize}[leftmargin=12pt, itemsep=0pt, topsep=1pt, parsep=0pt]
    \item \textbf{Certificación Excel Green Belt (A2 Capacitación):} Especialización en modelado analítico, funciones complejas y automatización.
    \item \textbf{Google Cloud Professional Database Engineer:} En preparación activa (arquitectura, disponibilidad y gestión en nube).
    \item \textbf{Especialización en Ciencia de Datos:} Formación continua en análisis exploratorio (EDA), manipulación de datos y estadística con Python.
\end{itemize}

\section{Competencias Profesionales}
\textbf{Enfoque Analítico:} Traducción de requerimientos operativos en soluciones de datos estructuradas y tableros accionables.\\
\textbf{Comunicación Técnica:} Habilidad para presentar análisis y hallazgos tanto a equipos de desarrollo como a personal directivo.\\
\textbf{Calidad y Gobernanza:} Validación sistemática de registros, estandarización de consultas y resguardo de la integridad del dato.

\end{document}
